\section{Introducción}
El desarrollo de software contemporáneo enfrenta retos significativos, principalmente derivados de la creciente complejidad, escalabilidad y mantenibilidad de los sistemas. Estos desafíos impactan directamente en la calidad del código y su evolución a lo largo del tiempo. En este sentido, \cite{vera2022arquitectura} afirma que, sin una arquitectura de software adecuada, el desarrollo pierde efectividad, ya que la falta de organización afecta negativamente el ciclo de vida del software.

Para abordar estas problemáticas, la comunidad técnica y académica ha promovido la adopción de paradigmas y estructuras de diseño que fomentan la separación de responsabilidades y la reutilización del código. Según \cite{gavilanez2022analisis}, afirma que el paradigma de la Programación Orientada a Objetos (POO) se apoya en patrones de diseño y arquitecturas específicas como el Modelo-Vista-Controlador (MVC), las cuales buscan imponer modularidad, coherencia y claridad en el desarrollo.

De manera complementaria, los equipos de desarrollo han incorporado metodologías ágiles, destacando especialmente Scrum como marco de gestión de proyectos. Según \cite{marco2022revision} y \cite{merchan2022comparacion}, esta metodología propone ciclos iterativos cortos que favorecen la mejora continua, la entrega constante de valor y una comunicación más eficiente entre los miembros del equipo.

En este contexto, el presente trabajo busca analizar y aplicar estas prácticas tanto a nivel técnico como metodológico, con el objetivo de demostrar que la integración de arquitecturas de software bien estructuradas junto con marcos de trabajo ágiles favorece la creación de sistemas más robustos, mantenibles y alineados con las necesidades del cliente.